\emph{Pour chacune des affirmations suivantes, indiquer si elle est juste ou fausse. Chaque réponse doit être justifiée. Une réponse non justifiée ne rapporte aucun point.}

\medskip

\textbf{Affirmation 1 :} Soit (E) l'équation différentielle : $y'- 2 y= - 6 x + 1$.

La fonction $f$ définie sur $\R$ par :\quad $f(x)=\e^{2 x}-6 x + 1$\quad est une solution de l'équation différentielle ( E ).

\medskip

\textbf{Affirmation 2 :} On considère la suite $\left(u_{n}\right)$ définie sur $\N$ par

\[ u_{n}=1 + \dfrac{3}{4} + \left(\dfrac{3}{4}\right)^{2} + \cdots + \left(\dfrac{3}{4}\right)^{n}	\]

La suite $\left(u_{n}\right)$ a pour limite $ + \infty$.

\medskip

\textbf{Affirmation 3 :} On considère la suite $\left(u_{n}\right)$ définie dans l'affirmation 2.

L'instruction \texttt{suite(50)} ci-dessous, écrite en langage Python, renvoie $u_{50}$.

\begin{center}
	\begin{ttfamily}
	\begin{tabularx}{8cm}{c|X|}\cline{2-2}
		1&def suite(k):\\
		2&\quad S=0\\
		3&\quad for i in range(k):\\
		4&\qquad S=S+(3/4)**k\\
		5&\quad return S\\ \cline{2-2}
	\end{tabularx}
\end{ttfamily}
\end{center}

\medskip

\textbf{Affirmation 4 :} Soit $a$ un réel et $f$ la fonction définie sur ]0; $ + \infty[$ par :

\[f(x) = a \ln (x) - 2x.\]

Soit $C$ la courbe représentative de la fonction $f$ dans un repère \Oij.

II existe une valeur de $a$ pour laquelle la tangente à $C$ au point d'abscisse 1 est parallèle à l'axe des abscisses.
